\section{Проектиране на библиотеката}
\label{sec:design}

Настоящата глава представя проектните решения, върху които стъпва библиотеката за обектно-релационно картографиране (Object-Relational Mapping, ORM). Излагането върви от теоретичната основа на статичното моделиране, през слоестата организация и логическите слоеве (обекти, заявки, конектори), до асинхронната подсистема, изградена върху корутини (coroutines). Целта е читателят да види архитектурата като система от типове, формални договори (contracts) и проверими ограничения, а не като отделни езикови похвати.

\subsection{Теоретична основа на обектно-релационното картографиране}

Достъпът до данни е една от областите, в които разминаването между езика на приложението и модела на съхранение поражда най-много структурни грешки. Когато заявката се описва като низ, който се проверява едва при изпълнение, неправилните имена на полета, несъвместимите типове и невалидните клаузи се откриват късно. Класическият подход на обектно-релационното картографиране~\cite{fowler2002} свързва обектния модел на приложението с персистентното съхранение, но традиционните реализации остават силно зависими от изпълнителни механизми: текстови заявки, кодогенерация, анотации или динамично описани схеми.

В контекста на C++ това положение е особено противоречиво. Езикът предоставя мощна шаблонна (template) система, изчислимост по време на компилация (compile-time) чрез \code{constexpr} и концепции (concepts), но в библиотеките за достъп до данни тези механизми обикновено не се използват в пълна степен. Настоящата разработка изследва доколко самият компилатор може да поеме ролята на първи валидатор на структурата на заявката и на съвместимостта със системата за съхранение (backend).

Тук терминът обектно-релационно картографиране не се разбира тясно като техника, ограничена до релационни бази данни. Той обозначава архитектурен подход, при който моделът на данните, формата на заявките и част от ограниченията върху операциите се описват чрез типовата система на езика, а не чрез конфигурация по време на изпълнение (runtime). Това позволява една и съща логическа структура да бъде пренесена към релационни, документни, ключ-стойност (key-value), графови и ширококолонни (wide-column) системи, без различията им да се скриват или подменят.

\begin{figure}[H]
\centering
\begin{tikzpicture}[node distance=1.5cm and 1.4cm]
    \node[ormaccent] (cpp) {\textbf{C++ логически модел}\\обектни типове, полета, връзки};
    \node[ormblock, below=of cpp] (ir) {\textbf{Типизирано представяне}\\полета, правила, заместители};
    \node[ormblock, below left=of ir] (rel) {Релационен\\модел};
    \node[ormblock, below=of ir] (doc) {Документен\\модел};
    \node[ormblock, below right=of ir] (kv) {Ключ-стойност};
    \node[ormblock, below=of doc, xshift=-3.2cm] (graph) {Графов\\модел};
    \node[ormblock, below=of doc, xshift=3.2cm] (wide) {Ширококолонен\\модел};

    \draw[ormline] (cpp) -- (ir);
    \draw[ormline] (ir) -- (rel);
    \draw[ormline] (ir) -- (doc);
    \draw[ormline] (ir) -- (kv);
    \draw[ormline] (ir) -- (graph);
    \draw[ormline] (ir) -- (wide);
\end{tikzpicture}
\caption{Таксономия на логическия модел, междинното представяне на заявката и семействата системи за съхранение}
\label{fig:storage_taxonomy}
\end{figure}

Фигура~\ref{fig:storage_taxonomy} подчертава централното теоретично допускане на библиотеката: общото между различните системи за съхранение не е техният синтаксис, а логическият модел на данните и операциите. Затова валидната абстракция трябва да бъде разположена над структурирания език за заявки (Structured Query Language, SQL), двоичното разширение на JSON (Binary JSON, BSON), Redis командите, Cypher и Cassandra Query Language (CQL). Тя трябва да работи на равнището на полета, предикати, проекции, връзки и допустими операции.

Съществената особеност на подхода е, че моделирането по време на компилация не означава отсъствие на данни по време на изпълнение. Стойностите на параметрите се подават именно при изпълнение, но структурата на заявката е фиксирана и проверена още преди това. Получава се ясно разграничение между \emph{логическа форма} и \emph{стойности по време на изпълнение}, което е ключово за тезата на дипломната работа, защото позволява едновременно статични гаранции и работа с реални системи за съхранение.

\subsection{Модели за съхранение}

Библиотеката работи над пет основни модела за съхранение. Релационният модел организира данните в таблици, редове и колони и предоставя декларативен език за заявки със ясно дефинирани ограничения на целостта. Документният модел работи с колекции от документи, чиито полета могат да бъдат вложени и да варират по структура. Ключ-стойност моделът осигурява бърз достъп около първичен идентификатор, но има силно ограничени възможности за заявки върху предикати. Графовият модел представя данните като върхове, ребра и свойства, превръщайки връзките в първокласни обекти. Ширококолонният модел, представен от Cassandra, използва разпределени таблици с разделящи (partition) и подреждащи (clustering) ключове, при които правилната заявка следва физическата схема на разпределение.

\begin{table}[H]
\centering
\small
\caption{Съпоставка между основните модели за съхранение}
\label{tab:storage_models}
\begin{tabularx}{\textwidth}{L{2.3cm}L{2.7cm}L{2.8cm}Y}
\toprule
\multirow[t]{1}{=}{\textbf{Модел}} & \textbf{Единица за съхранение} & \textbf{Типични връзки} & \textbf{Основно ограничение} \\
\midrule
\multirow[t]{1}{=}{Релационен} & таблица / ред & външен ключ, свързване (JOIN) & необходимост от съгласувана схема и силна типова дисциплина \\
\hdashline
\multirow[t]{1}{=}{Документен} & документ & вграждане или препратка & не всяка релационна операция е естествена или евтина \\
\hdashline
\multirow[t]{1}{=}{Ключ-стойност} & ключ + стойност & справка по ключ, хеш-полета & силно ограничени заявки по предикати извън ключовия достъп \\
\hdashline
\multirow[t]{1}{=}{Графов} & връх + ребро & връзки за обхождане & специфичен език за заявки и графова семантика \\
\hdashline
\multirow[t]{1}{=}{Ширококолонен} & разделен ред & ключово ориентирани зависимости & заявките трябва да следват логиката на разделящия ключ \\
\bottomrule
\end{tabularx}
\end{table}

От Таблица~\ref{tab:storage_models} се вижда, че унифицирана абстракция е възможна само на равнището на логическия модел, полето, връзката и базовите операции. Отвъд тази граница започват различията: свързване няма универсален смисъл в Redis; вграждането няма идентична цена в релационна таблица и в документна колекция; обхождане на граф не се свежда до проста селекция; ограниченията, наложени от разделящия ключ в Cassandra, не могат да бъдат прикрити като обикновени правила за филтриране. Архитектурно правилният подход не е тези различия да се скриват, а да се направят явни и формално контролирани.

\subsection{Архитектурни цели и слоеста организация}
\label{subsec:architecture_goals}

Архитектурата на библиотеката е проектирана при едновременното действие на няколко ограничения. Първо, публичният програмен интерфейс трябва да остане в рамките на идиоматичен C++, без отделен специализиран език (Domain-Specific Language, DSL) и без отражение (reflection) по време на изпълнение. Второ, структурата на заявката трябва да бъде представима като \code{constexpr} обект, така че значима част от валидирането да се извършва от компилатора. Трето, слоят на конекторите трябва да бъде разширяем без наследяване от общ виртуален интерфейс, тъй като такъв интерфейс би скрил различията между системите за съхранение и би изместил проверките към изпълнение. Четвърто, библиотеката трябва да поддържа релационни и нерелационни конектори без фалшиво уеднаквяване --- общият програмен интерфейс и ограниченията на конкретната система за съхранение трябва да съществуват едновременно и проверимо.

\begin{figure}[H]
\centering
\begin{tikzpicture}[node distance=1.3cm and 1.4cm]
    \node[ormaccent] (app) {\textbf{Приложен слой}\\обектни типове и съставяне на заявки};
    \node[ormblock, below=of app] (entity) {\textbf{Обектен слой}\\\code{property}, \code{relationship}, \code{table\_name\_trait}};
    \node[ormblock, below=of entity] (ir) {\textbf{Слой на заявките}\\\code{field}, \code{Rule}, \code{select\_query}, \code{insert\_query}};
    \node[ormblock, below=of ir] (connector) {\textbf{Слой на конекторите}\\\code{connector\_trait<DB>} и способности (capabilities)};
    \node[ormblock, below=of connector] (driver) {\textbf{Драйверен слой}\\материализация към SQL / BSON / Redis / Cypher / CQL};

    \draw[ormline] (app) -- (entity);
    \draw[ormline] (entity) -- (ir);
    \draw[ormline] (ir) -- (connector);
    \draw[ormline] (connector) -- (driver);
\end{tikzpicture}
\caption{Слоеста архитектура на библиотеката}
\label{fig:architecture_layers}
\end{figure}

Фигура~\ref{fig:architecture_layers} показва последователната вертикална организация. Приложният код не общува пряко с драйвера, а работи с C++ типове и съставители на заявки (query builders). Драйверът от своя страна не познава семантиката на областта, а получава вече материализираната форма, генерирана от слоя на конекторите. Когато потребителят смени \code{SQLiteDB} с \code{MongoDB}, междинното представяне на заявката остава логически същото, но конекторът го тълкува като изпълнение върху филтър и проекция вместо изобразяване към SQL. Това не означава, че всяко междинно представяне е допустимо върху всяка система за съхранение; означава, че \emph{формата} на абстракцията е обща, а допустимостта се определя формално чрез механизма на способностите.

\subsection{Обектен слой}
\label{subsec:entity_layer}

Обектният слой е изграден около три централни конструкции: \code{property<T, Name>}, \code{relationship<...>} и \code{table\_name\_trait<T>}. Те имат двойна роля --- представят C++ модела на областта и същевременно описват логическата структура на персистентните данни. Така в библиотеката не съществува отделен регистър от метаданни по време на изпълнение; типовата система съдържа достатъчно информация, за да бъдат извлечени имената на полетата, типовете, връзките и целевият контейнер за съхранение.

Имената на логическите контейнери чрез \code{table\_name\_trait<T>} играят обобщена роля. В SQLite, PostgreSQL и MySQL те означават таблица; в MongoDB --- колекция; в Redis --- ключов префикс; в Neo4j --- етикет (label). Единен логически интерфейс адресира физически различни структури, а авторът на конектора извлича схемата директно от типовете, а не от външен регистър.

\begin{lstlisting}[language=C++,caption={Извадка от \code{property} и \code{relationship} в обектния слой},label={lst:property_relationship}]
template <typename T, detail::string_literal Name>
struct property
{
    using value_type = T;

    [[nodiscard]] static constexpr std::string_view column_name() noexcept
    {
        return Name.view();
    }

    T value{};
};

enum class store_as { reference, embed };

template <store_as Strategy, typename T, detail::string_literal Name>
struct relationship
{
    static constexpr store_as strategy = Strategy;
    using related_type = T;
};
\end{lstlisting}

Listing~\ref{lst:property_relationship} показва, че обектният слой е минималистичен по форма, но достатъчно богат по смисъл. \code{property} не е просто контейнер за стойност, а носи и компилационно име на колоната. \code{relationship} не е обикновен член, а фиксира стратегията за материализация като част от типа: \code{store\_as::reference} обозначава логическа външна препратка (foreign key, справка, ребро в граф), а \code{store\_as::embed} обозначава физическо включване в родителската структура (вложен документ, денормализирани данни). Връзката се описва веднъж на логическо ниво, а интерпретацията остава задача на конкретния конектор.

\subsection{Слой на заявките}
\label{subsec:query_layer}

Вместо заявката да се описва пряко като SQL или друго текстово представяне, библиотеката използва типизирано междинно представяне на заявката (query intermediate representation, query IR). Неговите елементи са \code{field}, \code{Rule}, \code{Placeholder} и обектите \code{select\_query}, \code{insert\_query}, \code{update\_query} и \code{delete\_query}. Този модел има две предимства. Първо, структурата на заявката може да бъде анализирана по време на компилация и да участва в \code{static\_assert} проверки. Второ, едно и също междинно представяне може да бъде преведено към различни езици и протоколи за съхранение.

Концептуално междинното представяне играе ролята на междинен език между модела на областта и слоя за изпълнение. Това е аналогично на междинните представяния в компилаторите: потребителят работи с високо ниво на абстракция, а конкретната система за съхранение получава специализирана форма, оптимизирана за собствения си модел.

\begin{lstlisting}[language=C++,caption={Извадка от \code{select\_query} и съставяне на правила},label={lst:select_query_excerpt}]
template <typename Response, typename Joins, typename Wheres,
          typename Limits, typename Groups, typename Orders>
struct select_query : select_query_tag
{
    using response_type = Response;
    using row_type = projected_type<Response>;

    template <typename RuleType>
        requires is_rule<RuleType>
    [[nodiscard]] constexpr auto where(RuleType r) const
    {
        using NewWheres = detail::append_type_t<Wheres, RuleType>;
        return select_query<Response, Joins, NewWheres, Limits, Groups, Orders>(
            response_, joins_,
            detail::tuple_cat(wheres_, detail::orm_tuple<RuleType>(r)),
            limits_, groups_, orders_);
    }
};
\end{lstlisting}

Listing~\ref{lst:select_query_excerpt} илюстрира важно теоретично свойство: съставителят на заявки не променя низ, а изгражда нови типизирани обекти. Последователността \code{select(...).where(...).group\_by(...)} представлява конструиране на нов тип, а не серия от странични ефекти върху текстова форма.

Полето в израза на заявката се представя чрез \code{field<\&T::m>} --- това е \code{constexpr} обвивка над указател към член, която носи типа на контейнера, типа на стойността и името на колоната. Операторите върху \code{field} създават типизирани правила (\code{Rule}). Изразът \code{field<\&User::id> == Placeholder<int>\{\}} не е низ, а статично описание на предикат.

Механизмът на заместителите има две форми. \code{Placeholder<T>} моделира анонимни параметри, които се консумират позиционно от аргументите на \code{execute()}. \code{orm::ph<T, std::placeholders::\_N>} моделира индексирани параметри, при които една стойност по време на изпълнение може да се използва на повече от едно място. Тази възможност е особено полезна за системи за съхранение като SQLite, PostgreSQL и MySQL, при които позицията на заместителя е отделна част от логиката на материализация. Тъй като веригата на съставителя е съвместима с \code{constexpr}, значима част от формата на заявката може да бъде изградена и проверена още при компилация (Listing~\ref{lst:constexpr_query}).

\begin{lstlisting}[language=C++,caption={Заявка декларирана като константа, изчислима при компилация},label={lst:constexpr_query}]
using namespace orm::literals;
static constexpr auto get_active =
    orm::select(orm::field<&User::id>, orm::field<&User::name>)
        .where(orm::field<&User::score> > 0.0)
        .where(orm::field<&User::id> == orm::Placeholder<int>{})
        .limit(10_per_page & 2_page);
\end{lstlisting}

\subsection{Слой на конекторите и формален договор}
\label{subsec:connector_contract}

Конекторът се описва чрез тип \code{DB} и специализация \code{connector\_trait<DB>}. Това е структурен, а не наследствен договор: няма общ виртуален интерфейс \code{IConnector}, защото различните системи за съхранение имат различни възможности и ограничения, които трябва да бъдат видими по време на компилация. Виртуалната йерархия би изравнила твърде много различия и би прехвърлила част от проверките към изпълнение.

Минималният работещ конектор трябва да предостави \code{wire\_type}, \code{cursor\_type} и подходящи претоварвания на \code{execute(...)}. Допълнителните възможности се обозначават чрез маркери за способности (capability tags). По този начин договорът остава едновременно строг и разширяем; компилаторът вижда специализацията директно и може да провери дали тя удовлетворява концепцията \code{is\_connector}.

\begin{lstlisting}[language=C++,caption={Извадка от формалния договор на конектора},label={lst:connector_contract}]
template <typename DB>
concept is_connector = requires {
    typename connector_trait<DB>::template wire_type<int>::type;
    typename connector_trait<DB>::cursor_type;
};

template <typename DB, typename Cap>
inline constexpr bool has_capability =
    detail::capability_check<DB, Cap>::value;
\end{lstlisting}

Listing~\ref{lst:connector_contract} показва, че договорът е структурен, а не наследствен. Това е решаващ архитектурен избор: ако съществуваше виртуален интерфейс с диспечеризация по време на изпълнение, част от типово проверимата информация би била изгубена. Вместо това библиотеката предпочита статична структура, при която изискванията са видими както за компилатора, така и за автора на конектора.

\begin{table}[H]
\centering
\small
\caption{Задължителни и опционални елементи на конектора}
\label{tab:connector_requirements}
\begin{tabularx}{\textwidth}{L{5.5cm}Y}
\toprule
\textbf{Елемент} & \textbf{Роля} \\
\midrule
\multirow[t]{1}{=}{\code{DB} тип} & обвивка около връзка, сесия или имитиращ (mock) обект \\
\hdashline
\multirow[t]{1}{=}{\code{connector\_trait<DB>}} & централна специализация на договора \\
\hdashline
\multirow[t]{1}{=}{шаблон \code{wire\_type<T>}} & превръща C++ тип в представяне, използвано от драйвера (wire representation) \\
\hdashline
\multirow[t]{1}{=}{\code{cursor\_type}} & описва обхождането и жизнения цикъл на резултатите \\
\hdashline
\multirow[t]{1}{=}{\code{execute(...)}} & материализира междинното представяне към действие на системата за съхранение \\
\hdashline
\multirow[t]{1}{=}{Маркери за способности} & ограничават допустимите операции и политики на жизнения цикъл \\
\hdashline
\multirow[t]{1}{=}{\code{begin}/\code{commit}/\code{rollback}} & необходими при поддръжка на транзакции \\
\hdashline
\multirow[t]{1}{=}{\code{ddl\_for(...)}} & необходимо при интеграция с миграция на схемата \\
\bottomrule
\end{tabularx}
\end{table}

Имената на конекторите следват ясна конвенция: \code{SQLiteDB}, \code{PostgreSQLDB}, \code{MySQLDB}, \code{MongoDB}, \code{RedisDB}, \code{Neo4jDB} и \code{CassandraDB} обозначават локални или mock-подобни обвивки, докато реалните варианти, които ползват жив (live) драйвер, носят суфикс \code{LiveDB}. Това не е стилистичен избор, а част от яснотата на договора: още от името личи дали типът е тестов адаптер, локална обвивка или система за съхранение, обвързана с реален драйвер.

\subsection{Механизъм на способностите}
\label{subsec:capabilities}

Маркерите за способности са основният начин библиотеката да различава конекторите по допустими операции. Ако \code{connector\_trait<DB>} не декларира \code{supports\_joins}, опитът за заявка със свързване (JOIN) е архитектурно неправилен и трябва да бъде отхвърлен още при компилация. Същото важи за транзакциите, за агрегацията и за конкурентното изпълнение. По този начин коректността не зависи единствено от тестове --- част от договора се проверява пряко от компилатора.

\begin{table}[H]
\centering
\small
\caption{Основни маркери за способности и тяхното значение}
\label{tab:capability_tags}
\begin{tabularx}{\textwidth}{L{5.5cm}Y}
\toprule
\textbf{Способност} & \textbf{Архитектурно значение} \\
\midrule
\multirow[t]{1}{=}{\code{supports\_joins}} & разрешава заявки със свързване за дадената система за съхранение \\
\hdashline
\multirow[t]{1}{=}{\code{supports\_transactions}} & позволява \code{transaction\_guard} и куки за \code{begin}/\code{commit}/\code{rollback} \\
\hdashline
\multirow[t]{1}{=}{\code{supports\_aggregation}} & обозначава поддръжка на агрегиращи операции в слоя на заявките \\
\hdashline
\multirow[t]{1}{=}{\code{supports\_embedding}} & декларира естествена поддръжка на вградени структури \\
\hdashline
\multirow[t]{1}{=}{\code{supports\_upsert}} & маркира системи за съхранение, при които вмъкване-или-обновяване (upsert) е част от естествения договор \\
\hdashline
\multirow[t]{1}{=}{\code{supports\_bulk\_insert}} & обозначава наличие на ефективен път за масово записване \\
\hdashline
\multirow[t]{1}{=}{\code{supports\_concurrent\_execute}} & необходимо за \code{connection\_pool<DB, N>} \\
\hdashline
\multirow[t]{1}{=}{\code{supports\_async}} & задължително за интеграция чрез неблокиращите интерфейси на драйвера в асинхронния слой \\
\bottomrule
\end{tabularx}
\end{table}

Класът \code{db<DB>} е публичното фасадно ниво. Той обвива препратка към конкретна система за съхранение и предоставя унифициран интерфейс за изпълнение. Същевременно той е архитектурен филтър: точно тук статичните проверки за способности и съвместимост се срещат с конкретното междинно представяне на заявката.

\begin{lstlisting}[language=C++,caption={Проверка на възможностите във фасадата \code{db<DB>}},label={lst:db_facade}]
template <typename DB>
    requires is_connector<DB>
class db
{
public:
    template <typename Query>
    auto operator<<(Query q)
    {
        if constexpr (is_select_query<Query>)
        {
            if constexpr (decltype(q.join_clauses())::size > 0)
            {
                static_assert(has_capability<DB, cap::supports_joins>,
                    "orm::db: this connector does not support JOIN operations.");
            }
        }
        return connector_trait<DB>::execute(*conn_, std::move(q));
    }
};
\end{lstlisting}

Listing~\ref{lst:db_facade} показва защо \code{db<DB>} не е тънка обвивка, а формалното място, на което структурата на заявката и информацията за способностите на конектора се срещат. Именно тук статичната архитектура се превръща в реално ограничение върху употребата на интерфейса. Аналогична логика прилагат и помощниците \code{connection\_pool<DB, N>}, \code{thread\_local\_db<DB>} и \code{transaction\_guard<DB>}: всеки от тях изисква декларацията на съответната способност в \code{connector\_trait<DB>}, без която компилацията не преминава.

Особено важно е, че проверката на способностите включва и негативни декларации: ако MongoDB или Redis не поддържат свързване и транзакции, отсъствието на съответните маркери също е част от договора и се валидира от тестовете на конкретните конектори. Това различава нашата архитектура от наследствените йерархии, в които „наследниците“ са длъжни да реализират цялостен интерфейс независимо дали той има смисъл за тяхната семантика.

\subsection{Прототипен слой за миграции}

Архитектурата включва и прототипен слой за миграция на схемата чрез \code{migrate<DB>}, \code{live\_schema}, \code{entity\_meta} и операции на езика за описание на данни (Data Definition Language, DDL) като \code{create\_table\_op}, \code{add\_column\_op}, \code{drop\_column\_op} и \code{alter\_column\_type\_op}. Този слой не е универсално завършен за всяка система за съхранение, но архитектурно показва, че библиотеката не се ограничава до изпълнение на заявки. При достатъчно богата специализация на \code{connector\_trait<DB>} обектното описание участва и в развитието на схемата. Това променя смисъла на слоя за обектно-релационно картографиране: той се превръща не само в интерфейс към данните, но и в потенциален източник на инструменти, осведомени за схемата (schema-aware tooling). Конкретното изпълнение на тази подсистема и подкрепящите я тестове са описани в Глава~\ref{sec:implementation}.

\subsection{Теория на корутините и асинхронен дизайн}
\label{subsec:coroutines_design}

Втората основна архитектурна ос на разработката е асинхронното изпълнение. Корутините (coroutines) са езиков механизъм за описване на прекъсваеми изчисления, при които изпълнението може да бъде спряно и по-късно възобновено без загуба на локалното състояние. В стандарта C++20 те влизат в основния език~\cite{cpp20standard,p0912r5,p1063r0}, което позволява асинхронният слой да се изгражда върху статично дефинирани правила за типове, жизнен цикъл и поток на управление, вместо върху ad hoc обратни извиквания (callbacks).

\subsubsection{Корутини спрямо нишки и процеси}

Процесът е защитен адресен контекст на операционната система, нишката е планирана от ядрото последователност от инструкции, а корутината е кооперативна единица за изпълнение, чието спиране и възобновяване се управляват от програмния модел, а не пряко от планировчика на ядрото. Това разграничение е съществено, защото библиотеката комбинира корутини и обединение от нишки (thread pool) без да приравнява двете абстракции. Нишката е скъп ресурс на ядрото със собствен стек и синхронизационни примитиви; корутината е значително по-лека логическа единица, която пази само състоянието, необходимо за конкретното изчисление.

\begin{table}[H]
\centering
\small
\caption{Съпоставка между основните изпълнителни единици}
\label{tab:execution_units_comparison}
\begin{tabularx}{\textwidth}{L{2.8cm}L{4.0cm}Y}
\toprule
\textbf{Единица} & \textbf{Същност} & \textbf{Значение за библиотеката} \\
\midrule
Процес & Изолиран адресен контекст с ресурси на операционната система & Не е пряка единица на архитектурата, но определя средата за изпълнение \\
\hdashline
Нишка & Поток на изпълнение, планиран от ядрото & Използва се от обединението от нишки за изнасяне на блокиращи операции \\
\hdashline
Корутина & Кооперативно прекъсваемо изчисление със запазено локално състояние & Осигурява формалния модел за \code{Task<T>} и за асинхронната композиция \\
\bottomrule
\end{tabularx}
\end{table}

\subsubsection{Awaitable договор и преобразуване в крайно автоматно представяне}

Ключовите езикови оператори \code{co\_await}, \code{co\_return} и \code{co\_yield} маркират, че функцията следва корутинна семантика. Най-важен е \code{co\_await}: той описва точката, в която текущото изчисление отстъпва изпълнението, докато външна операция --- база данни или планиране към обединението от нишки --- стане готова. Така синтаксисът на потребителя остава линеен, докато моделът на изпълнение е асинхронен.

Обектът, който се изчаква, участва в тристъпков договор. \code{await\_ready()} определя дали изобщо е необходимо спиране; \code{await\_suspend()} получава манипулатор на корутината и решава как ще се организира възобновяването; \code{await\_resume()} връща резултата или изнася изключение. Тази тройка обединява разнородни източници на асинхронност --- задачи в обединение от нишки, готовност на канал в реактор, обратно извикване от драйвер --- под единен синтаксис на \code{co\_await}.

Компилаторът преобразува корутинната функция в крайно автоматно представяне (state machine) с ясно дефинирани точки на спиране, преходи на възобновяване и кадрово (frame) съхранение~\cite{p0912r5,p1063r0}. Това е особено значимо за тезата на работата: същият статичен инструментариум, който описва заявките, дава възможност и потокът на асинхронно управление да бъде формализиран чрез типове и статични правила, вместо чрез неструктурирани вериги от обратни извиквания.

\subsubsection{\code{Task<T>} като носител на жизнения цикъл}

Основната потребителска абстракция е \code{Task<T>}. Тя не се изчерпва с „връща бъдещ резултат“ --- тя е контейнер за корутинния кадър, обекта на обещанието (promise object), продължението (continuation) и политиката за завършване. Реализацията използва \code{initial\_suspend()} и \code{final\_suspend()} така, че корутината по подразбиране е мързелива (lazy) и предава управлението обратно към продължението при завършване. Така \code{Task<T>} е едновременно ръкохватка за изпълнение и формален договор за това как се предават резултат, изключение и завършване между корутинни контексти.

От инженерна гледна точка са важни три свойства. Първо, \code{operator co\_await()} свързва текущото продължение към изчакваната задача и позволява естествено вложено съставяне с \code{co\_await}. Второ, \code{sync\_wait()} осигурява мост към некорутинен код, което е необходимо за единичните (unit) тестове и за граничните сценарии, в които приложението още не работи изцяло в корутинен стил. Трето, \code{detach()} и \code{start\_detached()} позволяват контролирано поведение „пусни и забрави“, при което кадърът на корутината се освобождава самостоятелно след \code{final\_suspend()}.

\subsubsection{Универсален и асинхронен път през интерфейса на драйвера}

Над тази основа стои \code{async\_db<DB>} --- асинхронната фасада, която обвива \code{db<DB>} и запазва синтаксиса на заявката, но връща \code{Task<result<...>>}. Архитектурно тук се случва най-важният избор: \code{operator<<} съдържа \code{if constexpr} разклонение върху \code{has\_capability<DB, cap::supports\_async>}. Ако конекторът декларира, че драйверът предоставя асинхронен интерфейс, се извиква \code{connector\_trait<DB>::async\_execute()}. В противен случай синхронното изпълнение се обвива в \code{run\_on\_pool()}, който публикува заявката към \code{ThreadPool}, спира извикващата корутина и я възобновява, когато работата приключи.

\begin{lstlisting}[language=C++,caption={Разклонение на логиката в зависимост от възможностите на съответния конектор в \code{async\_db<DB>}},label={lst:async_db_dispatch}]
template <typename Query>
[[nodiscard]] auto operator<<(Query q)
    -> Task<decltype(std::declval<db<DB>>() << std::declval<Query>())>
{
    using ResultType = decltype(std::declval<db<DB>>() << std::declval<Query>());

    if constexpr (has_capability<DB, cap::supports_async>)
    {
        co_return co_await connector_trait<DB>::async_execute(
            db_.connection(), std::move(q));
    }
    else
    {
        co_return co_await run_on_pool(
            *pool_, [this, q = std::move(q)]() mutable -> ResultType {
                return db_ << std::move(q);
            });
    }
}
\end{lstlisting}

Listing~\ref{lst:async_db_dispatch} е централен за цялата асинхронна архитектура. Той показва, че библиотеката не използва регистър по време на изпълнение или диспечеризация на основата на низове, а статична проверка върху способностите на конектора. По този начин потребителският интерфейс остава еднороден, а вътрешният път на изпълнение остава честен спрямо реалните възможности на конкретната система за съхранение.

\begin{figure}[H]
\centering
\begin{tikzpicture}[node distance=1.4cm and 1.6cm]
    \node[ormaccent, align=center] (call) {Потребителска корутина\\\code{co\_await (adb << query)}};
    \node[ormblock, below=of call, align=center] (adb) {\code{async\_db<DB>}};
    \node[ormblock, below=of adb, align=center] (gate) {Статично разклонение\\\code{has\_capability<DB, cap::supports\_async>}};

    \node[ormblock, below left=of gate, align=center] (poolawait) {\code{run\_on\_pool()}};
    \node[ormblock, below=of poolawait, align=center] (syncpath) {\code{ThreadPool} +\\\code{db<DB>::execute}};

    \node[ormblock, below right=of gate, align=center] (native) {\code{connector\_trait<DB>::}\\\code{async\_execute}};
    \node[ormblock, below=of native, align=center] (iopath) {\code{IoContext} +\\крайно автоматно представяне на драйвера};

    \node[ormaccent, below=of $(syncpath)!0.5!(iopath)$, align=center] (result) {\code{Task<result<...>>}\\продължение и резултат};

    \draw[ormline] (call) -- (adb);
    \draw[ormline] (adb) -- (gate);
    \draw[ormline] (gate) -- (poolawait);
    \draw[ormline] (poolawait) -- (syncpath);
    \draw[ormline] (gate) -- (native);
    \draw[ormline] (native) -- (iopath);
    \draw[ormline] (syncpath) -- (result);
    \draw[ormline] (iopath) -- (result);
\end{tikzpicture}
\caption{Двата пътя на изпълнение на \code{async\_db<DB>}: универсално изнасяне и интеграция през асинхронния интерфейс на драйвера}
\label{fig:async_dispatch_paths}
\end{figure}

За драйверите, които предоставят неблокиращи интерфейси, библиотеката не се ограничава до изнасяне в обединение от нишки. Тук влиза в действие \code{IoContext} --- абстракция за реакторен цикъл, която предоставя \code{run()}, \code{post()}, \code{watch\_readable(fd)} и \code{watch\_writable(fd)}. Тези обекти не извършват самостоятелно входно-изходна операция; те само спират корутината и я възобновяват, когато съответният файлов описател стане готов за четене или запис. Това разделение между готовност и действие е решаващо: \code{IoContext} не „крие“ драйвера, а предоставя реакторен интерфейс, върху който клиентската библиотека с неблокиращ интерфейс може да стъпи (например \code{PQconnectPoll()}, \code{PQsendQueryParams()}, \code{PQflush()} и \code{PQisBusy()} на PostgreSQL).

Така асинхронният път през интерфейса на драйвера не е „по-бързо обединение от нишки“, а качествено различен модел: работната нишка не блокира в очакване на мрежова готовност; корутината се спира, а възобновяването се управлява от наблюдението на реактора. В Глава~\ref{sec:implementation} този модел се проявява конкретно при PostgreSQL, докато за SQLite, MySQL, MongoDB и Redis библиотеката използва изнасяне в обединение от нишки като честен и адекватен заместител, защото съответните драйвери нямат неблокираща семантика със същата зрелост.

\subsubsection{Координация на връзките и транзакциите при асинхронен достъп}

Асинхронната архитектура не приключва с единичното изпълнение на заявка. За реално приложение е необходимо управление на връзките и транзакциите при конкурентен достъп. За тази цел библиотеката предоставя \code{async\_connection\_pool<DB, N>}, \code{async\_connection\_guard<DB>}, \code{async\_transaction\_guard<DB>} и фабричната корутина \code{async\_begin\_transaction()}. \code{async\_connection\_pool<DB, N>} е продължение на синхронния обединителен слой, интегрирано с корутинния модел: методът \code{acquire()} връща \code{Task<async\_connection\_guard<DB>>} и вътрешно използва \code{run\_on\_pool()} и условна променлива, за да изчака свободен слот, без извикващата корутина да бъде изложена на блокиращ интерфейс.

Същата логика се наблюдава и при транзакциите. \code{async\_begin\_transaction()} първо изнася \code{BEGIN} в обединението от нишки и след това връща \code{async\_transaction\_guard<DB>}. Методите \code{commit()} и \code{rollback()} са \code{Task<void>} операции, а деструкторът на пазача (guard) изпълнява защитно \code{ROLLBACK}, ако транзакцията не е финализирана изрично. Дизайнът е едновременно удобен и консервативен: успешният път е интегриран с корутинния модел, а пътят за почистване остава детерминистичен.

\subsection{Архитектурни инварианти и обобщение}

От разгледаните слоеве могат да бъдат формулирани четири инварианта. Първо, приложният код никога не зависи пряко от драйверен интерфейс. Второ, логиката на заявката се формулира в типизирано междинно представяне, а не в специфични за системата за съхранение низове. Трето, разширяването към нова система за съхранение се извършва чрез специализация на признаците (trait specialisation) и декларация на способности, а не чрез промяна на самия интерфейс за заявки. Четвърто, ако дадена операция е недопустима за избраната система за съхранение, това е видимо възможно най-рано --- за предпочитане при компилация.

Тези инварианти правят архитектурата едновременно инженерно подредена и академично защитима. Те позволяват всяко следващо разширение да бъде оценявано не само по това дали „работи“, а и по това дали запазва вече формулираните граници. Следващата глава представя как тези проектни решения се материализират в конкретен код: реализацията на асинхронния слой, конкретните конектори за SQLite, PostgreSQL, MySQL, MongoDB и Redis, и подробната верификация на цялата система.
